\chapter{Rotation Geometry}
\label{ch:ch08-rotation-geometry}

\section{Combining Rotations}

The previous chapter recovered angles from trigonometric ratios. A
different kind of question now presents itself, concerning not a single
angle but several angles acting in sequence. Suppose a point is rotated by
an angle $A$, and the result is then rotated again by an angle $B$. What is
the overall effect of the two rotations together? Does the outcome depend
on the order in which they are performed, and can two successive rotations
always be replaced by a single rotation? Questions of this kind move the
discussion from individual angles to transformations acting on the plane,
and, more importantly, they reveal the geometric origin of many identities
that are too often presented as isolated formulas to be memorized. The goal
of this chapter is to derive those identities from the geometry of rotation
itself.

\section{Rotations as Transformations}

Recall from Chapter~4 that every point on the unit circle can be written as
$(\cos\theta, \sin\theta)$, obtained by rotating the point $(1,0)$ through
angle $\theta$. Consider now an arbitrary point $(x,y)$, written in polar
form as $x = r\cos\phi$ and $y = r\sin\phi$. Rotating this point by an
angle $\theta$ changes only its angle, not its distance $r$ from the
origin, so the rotated point has coordinates $(r\cos(\phi+\theta),
r\sin(\phi+\theta))$. The task of this section is to express these new
coordinates in terms of $x$, $y$, and $\theta$ alone, without reference to
$\phi$ or $r$.

\begin{theorem}
Rotating a point $(x,y)$ counterclockwise through an angle $\theta$
produces the point
\[
(x',y') = (x\cos\theta - y\sin\theta,\ x\sin\theta + y\cos\theta).
\]
\end{theorem}

\begin{proof}
A rotation of the plane is a linear transformation fixing the origin: it
carries sums of displacements to sums of the rotated displacements, and it
preserves distance from the origin. Such a transformation is completely
determined by where it sends the two points $(1,0)$ and $(0,1)$, since
every point $(x,y)$ is the combination $x(1,0) + y(0,1)$. By the
definition of sine and cosine in Chapter~4, rotating $(1,0)$, which lies at
angle $0$, by $\theta$ sends it to the point at angle $\theta$, namely
$(\cos\theta, \sin\theta)$. The point $(0,1)$ lies a quarter turn ahead of
$(1,0)$, at angle $\pi/2$; a direct check across the four quadrants using
the special angles of Chapter~4 shows that a quarter turn always sends a
point $(a,b)$ on the unit circle to $(-b,a)$, so rotating $(0,1)$ by
$\theta$ sends it to $(-\sin\theta, \cos\theta)$. Combining the two images
by linearity,
\[
(x,y) = x(1,0) + y(0,1) \ \longmapsto\ x(\cos\theta,\sin\theta) +
y(-\sin\theta,\cos\theta) = (x\cos\theta - y\sin\theta,\ x\sin\theta +
y\cos\theta).
\]
\end{proof}

\section{Composition of Rotations}

Rotations obey a strikingly simple composition law.

\begin{theorem}
A rotation by angle $A$ followed by a rotation by angle $B$ is equivalent
to a single rotation by angle $A+B$.
\end{theorem}

\begin{proof}
Let a point initially have angular coordinate $\phi$. After the first
rotation its angular coordinate becomes $\phi + A$, and after the second
it becomes $(\phi + A) + B$. By associativity of addition, $(\phi+A)+B =
\phi + (A+B)$, so the final point is exactly the result of rotating once by
angle $A+B$.
\end{proof}

This simple fact is the source of the most important identities in
trigonometry.

\section{Deriving the Sum Formulas}

Consider the point $(\cos A, \sin A)$, obtained geometrically by rotating
$(1,0)$ through angle $A$. Rotating this point by an additional angle $B$
and applying the rotation formula of Section~8.2 gives
\[
x' = \cos A\cos B - \sin A\sin B, \qquad y' = \sin A\cos B + \cos A\sin B.
\]
By the composition theorem just proved, however, rotating first by $A$ and
then by $B$ is equivalent to rotating $(1,0)$ once by $A+B$, so the
resulting point must also equal $(\cos(A+B), \sin(A+B))$. Since both
descriptions refer to the same point, corresponding coordinates must agree:
\[
\boxed{\cos(A+B) = \cos A\cos B - \sin A\sin B}, \qquad
\boxed{\sin(A+B) = \sin A\cos B + \cos A\sin B}.
\]
These formulas are not arbitrary identities to be verified after the fact;
they are the coordinate expression of rotation composition, true because
two rotations performed in sequence must land on exactly the point that a
single combined rotation would reach.

\section{Difference Formulas}

The difference formulas follow immediately from the sum formulas, without
any further geometry. Replacing $B$ with $-B$ in the sum formulas and using
the symmetries proved in Chapter~6, $\cos(-B) = \cos B$ and $\sin(-B) =
-\sin B$, gives
\[
\boxed{\cos(A-B) = \cos A\cos B + \sin A\sin B}, \qquad
\boxed{\sin(A-B) = \sin A\cos B - \cos A\sin B}.
\]
Nothing new has been assumed here; the difference formulas are simply the
sum formulas viewed through the even symmetry of cosine and the odd
symmetry of sine established earlier in the book.

\begin{example}
Compute $\sin(75^\circ)$ exactly. Writing $75^\circ = 45^\circ + 30^\circ$
and applying the sum formula for sine,
\[
\sin(75^\circ) = \sin45^\circ\cos30^\circ + \cos45^\circ\sin30^\circ
= \left(\frac{\sqrt2}{2}\right)\left(\frac{\sqrt3}{2}\right) +
\left(\frac{\sqrt2}{2}\right)\left(\frac12\right)
= \frac{\sqrt6+\sqrt2}{4}.
\]
\end{example}

\begin{algorithmproc}{How to Derive an Identity from Rotation Composition}
Identify a point obtained as the image of a known point under one
rotation, then apply a second rotation to it using the rotation formula of
Section~8.2 to obtain its coordinates algebraically. Separately, express
the same final point using angle addition and the composition theorem.
Since both expressions describe the identical point, equate corresponding
coordinates and simplify the resulting equations; the result is the
desired identity. Many of the identities appearing later in this book arise
from precisely this strategy.
\end{algorithmproc}

\begin{practicenow}
Compute $\cos(15^\circ)$ exactly, using $15^\circ = 45^\circ - 30^\circ$.
\end{practicenow}

\begin{sidebar}{Arabic Root-and-Pattern Morphology}
Arabic and the other Semitic languages build words through a mechanism
unusually close to the operator-and-object structure of this chapter. A
word is generated from a consonantal root, typically three consonants
carrying a core meaning, together with a separately specified pattern, a
template of vowels and affixes into which the root's consonants are
inserted. The root $f$-$\text{'}$-$l$, associated broadly with the idea
of doing or acting, yields distinct words such as \textit{fa'ala} (he
did), \textit{fa\text{''}ala} (he caused to do, intensively),
\textit{f\={a}'ala} (he did reciprocally, with another), and
\textit{af'ala} (he caused to do), among others, depending only on which
pattern is applied to the same fixed root. The root supplies the
invariant content; the pattern supplies the transformation acting on it,
and in each case the output can be described as
$\text{Pattern}(\text{Root}) = \text{Word}$, a fixed object transformed
by a chosen operator into a new but recognizably related object. This is
precisely the structure a rotation exhibits in this chapter, where a
fixed point is transformed by an operator $R(\theta)$ into a new point
that is different in position yet identical in its distance from the
origin: an invariant core, and a family of operators acting on it to
produce variation without destroying what makes each output recognizable
as a transformation of the same underlying thing.
\end{sidebar}

\section{Rotation Matrices}

The rotation transformation of Section~8.2 can be written compactly as
\[
\begin{pmatrix} x' \\ y' \end{pmatrix}
=
\begin{pmatrix} \cos\theta & -\sin\theta \\ \sin\theta & \cos\theta \end{pmatrix}
\begin{pmatrix} x \\ y \end{pmatrix}.
\]
This array of numbers is called a matrix, and although a full study of
matrices lies beyond the scope of this book, one fact is worth noting here:
composing two rotations, as in Section~8.4, corresponds exactly to
multiplying their two matrices together. The matrix form therefore encodes
the same angle-addition structure derived earlier in this chapter, in a
notation that will reappear in later mathematical study, particularly in
linear algebra, physics, and computer graphics.

\section{Applications}

Rotations occur throughout science and technology: rotating objects in
computer graphics, steering the joints of a robotic arm, controlling the
orientation of a spacecraft, describing crystallographic symmetries, and
analyzing mechanical linkages and gears all depend on the transformation
studied in this chapter. The identities derived above are not merely
algebraic conveniences to be looked up when needed; they are mathematical
descriptions of how rotational transformations combine, and every one of
these applications relies on that combination rule holding exactly.

\section{Looking Ahead}

The sum and difference formulas mark a turning point in the book. A large
share of the identities encountered in the chapters that follow are direct
consequences of them: the double-angle formulas of Chapter~10 arise simply
by setting $A = B$ in the sum formulas derived here, the half-angle
formulas emerge by solving those same equations in reverse, and the
product-to-sum formulas follow from further algebraic manipulation of the
same starting point. A large portion of trigonometric algebra therefore
grows from a single geometric fact established in this chapter: rotations
compose by adding angles. The next chapters develop this identity toolkit
further, showing how increasingly complex trigonometric expressions can be
simplified, transformed, and ultimately solved.

\section{Exercises}
\begin{exercise}
Use the sum formula to compute $\sin(105^\circ)$ exactly.
\end{exercise}

\begin{exercise}
Use the difference formula to compute $\cos(75^\circ)$ exactly.
\end{exercise}

\begin{exercise}
Use the sum formula to compute $\sin(15^\circ)$ exactly.
\end{exercise}

\begin{exercise}
Use the difference formula to compute $\cos(105^\circ)$ exactly.
\end{exercise}

\begin{exercise}
Starting from the sum formula for $\sin(A+B)$, derive the difference
formula for sine.
\end{exercise}

\begin{exercise}
Starting from the sum formula for $\cos(A+B)$, derive the difference
formula for cosine.
\end{exercise}

\begin{exercise}
Verify directly that two successive rotations of $40^\circ$ and $50^\circ$
produce the same result as a single rotation of $90^\circ$.
\end{exercise}

\begin{exercise}
Explain why a rotation cannot change the distance of a point from the
origin.
\end{exercise}

\begin{exercise}
Write the rotation matrix corresponding to a rotation of $90^\circ$, and
apply it to the point $(1,0)$.
\end{exercise}

\begin{exercise}
Why must composing a rotation by $A$ with a rotation by $B$ produce the
same result as a single rotation by $A+B$? What would break in the
geometry of the circle if this were not true?
\end{exercise}
